\begin{definition}[Base and integral marking]
\hypertarget{tag-01O1}{}\label{def:common-conventions}
\StatementRecord{def:common-conventions}{proved}{01O1}{}
Let $B=\mathbb P^1_t$, $p_1=0$, $p_2=1$, $p_0=\infty$,
$\Sigma=\{p_0,p_1,p_2\}$, and $B^\circ=B\setminus\Sigma$.
The full fibres of the marked family are $3S_1$, $4S_2$ and $W$;
$S_1,S_2,W$ are reduced, and $D=W+S_1+S_2$.
We orient base meridians counterclockwise in holomorphic local
coordinates and use clockwise root deck generators
$s_j(g_jz)=e^{-2\pi i/m_j}s_j(z)$, where $(m_1,m_2)=(3,4)$ and
$g_1g_2g_0=1$. The lift of a positive base loop ends at the inverse
deck translate, and its endpoint identification applies $g_j$.

The integral homology lattice and its dual basis are
\[
 \Lambda=\mathbb Z\langle\widehat\gamma,\widehat u,
                     \widehat w,\widehat\delta\rangle,
 \qquad (\gamma,u,w,\delta).
\]
The period matrix in coordinates $(\zeta_1,\zeta_2)$ is
\begin{equation}\hypertarget{tag-03CN}{}\label{wb:pshift:eq:periods}
 \Pi_c=\begin{pmatrix}6\mu&\tau&1&0\\
                    \beta_c&\mu&0&1\end{pmatrix},
 \qquad \beta_c=\beta^{\rm part}+c,
 \qquad (\beta^{\rm part}+\tau)(p_0)=0.
\end{equation}
Its rows are $\sigma_1=6\mu\gamma+\tau u+w$ and
$\sigma_2=\beta_c\gamma+\mu u+\delta$. The monodromy matrices are
\begin{equation}\hypertarget{tag-03C8}{}\label{wb:pshift:eq:finite-data}
 A_1=\begin{pmatrix}1&0&0&0\\6&0&1&0\\-6&-1&-1&0\\-2&1&0&1\end{pmatrix},
 \quad A_2=\begin{pmatrix}1&0&0&0\\0&0&-1&0\\-6&1&0&0\\3&0&1&1\end{pmatrix},
 \quad A_0=(A_1A_2)^{-1}.
\end{equation}
They act on columns in the displayed integral basis. Homology has
integral coefficients unless specified, and complex manifolds carry
their complex orientations. Reversing the two additive coordinates
acts by $J=\left(\begin{smallmatrix}0&1\\1&0\end{smallmatrix}\right)$:
the period and Hermitian matrices become $J\Pi_c$ and
$JHJ^{\mathsf t}$, respectively.
\end{definition}